% Leere Vorlage für Overleaf; LuaLaTeX oder pdfLaTeX.
\documentclass[11pt,a4paper]{article}
\usepackage[margin=22mm]{geometry}
\usepackage[T1]{fontenc}
\usepackage[utf8]{inputenc}
\usepackage[provide=*,ngerman]{babel}
\usepackage{tabularx,booktabs}
\usepackage[hidelinks]{hyperref}
\setlength{\parindent}{0pt}
\setlength{\parskip}{6pt}
\begin{document}
\section*{Sicherheitsprüfung: Berichtsvorlage}
\textbf{Vorlage, kein durchgeführter Prüflauf. Vertrauliche Rohdaten bleiben lokal.}
\begin{tabularx}{\linewidth}{lX}
\toprule
Berichts-ID / Gerätekennung & Ausfüllen \\
Beginn / Ende / Zeitzone & Ausfüllen \\
Ausführungsumgebung & Tatsächliches Gerät oder Cloud-Projekttest \\
Betriebssystem / Build & Ausfüllen \\
Codeversion / Commit / Hash & Ausfüllen \\
Werkzeuge / Erkennungsstand & Ausfüllen oder nicht eingesetzt \\
Vergleichsbericht & Ausfüllen oder keiner \\
\bottomrule
\end{tabularx}
\section*{Umfang und Veränderungen}
Erfasste Bereiche, Methoden und Vergleichbarkeit erläutern. Neue, veränderte,
entfernte und diesmal nicht erfasste Objekte unterscheiden.
\section*{Befunde und Unsicherheit}
Beleg, Einstufung, Unsicherheit und nächsten Schritt je Befund dokumentieren.
Kategorien: keine Hinweise im geprüften Umfang; unklar; schädliche Datei;
aktive Kompromittierung; nicht geprüft. Ein Inventar allein erhält kein
Sicherheitsurteil.
\section*{Lücken und Fehler}
Fehlende Rechte, unzugängliche Medien, verschlüsselte/übersprungene Dateien,
Zeit- und Größenlimits sowie weitere unberücksichtigte Bereiche nennen.
\section*{Maßnahmen}
Vorschläge von tatsächlich durchgeführten Maßnahmen unterscheiden.
\textbf{Datenverlust-Risiko:} Auswirkungen und Wiederherstellung vor destruktiven
Maßnahmen erläutern.
\section*{Abschluss}
Offene Aufgaben, nächste Prüfung, fachlich bestätigte Vergleichsbasis und
Datenschutzprüfung dokumentieren. Keine Garantie vollständiger Schadsoftwarefreiheit.
\end{document}
